% ============================================================
% 卒業論文アブストラクト本文（1頁）
% PDF上では自動的に2段組へ流れます．\switchcolumnは使いません．
% 各コメントは提出前に残っていてもPDFには出ませんが，内容を確認して
% 不要な節は整理してください．架空の数値や未実施の結果は書かないこと．
% ここには本文と\citeだけを書き，\printbibliographyや\nociteは置かないこと．
% main.texは埋込みアブスト用の一覧をアブスト内へ，本文用の一覧を巻末へ出力し，
% abstract.texは独立アブスト用の一覧を末尾へ出力します．
% ============================================================

\AbstractSection{1．はじめに}
% 研究背景，未解決課題，社会的・技術的な必要性を短く記載する．
% 最後に，本研究の目的を一文で明示する．

\AbstractSection{2．関連研究}
% 既存手法を列挙するのではなく，本研究との差分が分かるように整理する．
% 引用する場合はreferences.bibへ登録し，本文中で\cite{key}を使う．

\AbstractSection{3．提案・実施内容}
% 入力，出力，主要な処理，工夫点を，図や式も用いて具体的に記載する．
% 読み手が「何を新しく行ったか」を再現可能な粒度で示す．
% 図と番号付き数式を置いた場合は，本文中で図\ref{fig:...}，式\eqref{eq:...}のように参照する．
% 図題は図の下へ置く．

\AbstractSection{4．実験・評価}
% データ，分割，ベースライン，評価指標，実装条件，主要結果を記載する．
% 比較条件が同一であること，数値の単位とばらつきも確認する．
% 表を置いた場合は本文中で表\ref{tab:...}のように参照し，表題は表の上へ置く．

\AbstractSection{5．考察}
% 結果から直接いえることと推測を分け，成功例・失敗例・限界を述べる．

\AbstractSection{6．おわりに}
% 目的に対する到達点，得られた知見，残された課題を簡潔にまとめる．
